\subsection{The Three System Levels: SYS 0, SYS 1, SYS 2}
\label{subsec:11-3-system-levels}

%%%%%%%%%%%%%%%%%%%%%%%%%%%%%%%%%%%%%%%%%%%%%%%%%%%%%%%%%%%%%%%%%%%%%%%%%%%%%%%
% SECTION 11.3: The Three System Levels
% Module: BCM2_05_AWX_7240
% Content: SYS 0 (Visceral), SYS 1 (Affective), SYS 2 (Reflective)
%%%%%%%%%%%%%%%%%%%%%%%%%%%%%%%%%%%%%%%%%%%%%%%%%%%%%%%%%%%%%%%%%%%%%%%%%%%%%%%

Awareness operates across three distinct system levels, following the tradition
of Loewenstein's visceral factors and Kahneman's dual-process theory, extended
to a \textbf{triple-process model}.

%------------------------------------------------------------------------------
\subsubsection{Overview: Three Systems}
%------------------------------------------------------------------------------

\begin{table}[htbp]
\centering
\caption{Three System Levels of Awareness}
\label{tab:system-levels-overview}
\renewcommand{\arraystretch}{1.4}
\begin{tabular}{@{}lllll@{}}
\toprule
\textbf{Level} & \textbf{Name} & \textbf{Definition} & \textbf{Speed} & \textbf{Effort} \\
\midrule
SYS 0 & Instant/Visceral & Automatic activation without conscious processing & ms & None \\
SYS 1 & Affective/Emotional & Emotional awareness of meaning and resonance & sec & Low \\
SYS 2 & Reflective/Cognitive & Cognitive awareness of consequences and rules & min & High \\
\bottomrule
\end{tabular}
\end{table}

\paragraph{Extension Beyond Kahneman.}
While Kahneman's dual-process theory distinguishes System 1 (fast) from 
System 2 (slow), EBF adds a third level:

\begin{equation}
\text{Kahneman: } \{S1, S2\} \quad \rightarrow \quad \text{EBF: } \{SYS0, SYS1, SYS2\}
\end{equation}

SYS 0 captures \emph{visceral} processes that are even faster and more 
automatic than Kahneman's System 1.

%------------------------------------------------------------------------------
\subsubsection{SYS 0: Instant/Visceral Awareness}
%------------------------------------------------------------------------------

\begin{definition}[SYS 0: Instant Awareness]
\label{def:sys0}
Automatic visceral activation without conscious processing.
\end{definition}

\paragraph{Formal Condition.}
\begin{equation}
A_{SYS0} = 1 \quad \text{upon immediate sensory or affective reaction}
\end{equation}

\paragraph{Characteristics.}
\begin{itemize}
\item No cognitive mediation required
\item Evolutionarily hardwired responses
\item Cannot be suppressed (only regulated after the fact)
\item Dominant over SYS 1 and SYS 2 when activated
\item Operates in milliseconds
\end{itemize}

\paragraph{Examples.}
\begin{table}[htbp]
\centering
\caption{SYS 0 Awareness Examples}
\label{tab:sys0-examples}
\renewcommand{\arraystretch}{1.3}
\begin{tabular}{@{}lll@{}}
\toprule
\textbf{Trigger} & \textbf{Awareness Content} & \textbf{Utility Category} \\
\midrule
Pain stimulus & Immediate threat awareness & $U^{INU}_{P}$ \\
Hunger & Nutritional need awareness & $U^{INU}_{P}$ \\
Sexual arousal & Opportunity/threat awareness & $U^{INU}_{E}$ \\
Startle reflex & Danger awareness & $U^{INU}_{P}$ \\
Disgust reaction & Contamination awareness & $U^{INU}_{P}$ \\
Dopamine hit & Reward awareness & $U^{INU}_{E}$ \\
\bottomrule
\end{tabular}
\end{table}

\paragraph{Relation to Loewenstein's Visceral Factors.}
SYS 0 corresponds directly to Loewenstein's (1996, 2000) visceral factors:
drive states (hunger, thirst, sexual desire), emotions, and physical pain
that ``grab behavior'' automatically.

%------------------------------------------------------------------------------
\subsubsection{SYS 1: Affective/Emotional Awareness}
%------------------------------------------------------------------------------

\begin{definition}[SYS 1: Affective Awareness]
\label{def:sys1}
Emotional awareness of meaning and social resonance of one's own behavior.
\end{definition}

\paragraph{Formal Condition.}
\begin{equation}
A_{SYS1} = f(Emotion_{valence}, Social_{feedback})
\end{equation}

\paragraph{Characteristics.}
\begin{itemize}
\item Rapid but not instantaneous (seconds)
\item Socially calibrated
\item Can be influenced by context and framing
\item Interacts with both SYS 0 and SYS 2
\item Involves emotional processing
\end{itemize}

\paragraph{Examples.}
\begin{table}[htbp]
\centering
\caption{SYS 1 Awareness Examples}
\label{tab:sys1-examples}
\renewcommand{\arraystretch}{1.3}
\begin{tabular}{@{}lll@{}}
\toprule
\textbf{Emotion} & \textbf{Awareness Content} & \textbf{Utility Category} \\
\midrule
Guilt & Awareness of harm caused to others & $A^{comm}_{INU}$ \\
Pride & Awareness of achievement and status & $A^{self}_{IDN}$ \\
Shame & Awareness of norm violation/exposure & $A^{self}_{IDN}$ \\
Compassion & Awareness of others' suffering & $A^{comm}_{INU}$ \\
Envy & Awareness of relative deprivation & $A^{self}_{INU}$ \\
Belonging & Awareness of group connection & $A^{self}_{KNU}$ \\
\bottomrule
\end{tabular}
\end{table}

\paragraph{Social Calibration.}
SYS 1 awareness is inherently social---emotions like guilt, shame, pride,
and envy are calibrated relative to others and social norms.

%------------------------------------------------------------------------------
\subsubsection{SYS 2: Reflective/Cognitive Awareness}
%------------------------------------------------------------------------------

\begin{definition}[SYS 2: Reflective Awareness]
\label{def:sys2}
Cognitive awareness of consequences, rules, and normative implications.
\end{definition}

\paragraph{Formal Condition.}
\begin{equation}
A_{SYS2} = f(Causal_{understanding}, Rule_{clarity}, Evidence_{clarity})
\end{equation}

\paragraph{Characteristics.}
\begin{itemize}
\item Slow and effortful (minutes)
\item Requires cognitive capacity
\item Vulnerable to depletion and cognitive load
\item Can override SYS 0 and SYS 1 (with effort)
\item Involves deliberate reasoning
\end{itemize}

\paragraph{Examples.}
\begin{table}[htbp]
\centering
\caption{SYS 2 Awareness Examples}
\label{tab:sys2-examples}
\renewcommand{\arraystretch}{1.3}
\begin{tabular}{@{}lll@{}}
\toprule
\textbf{Process} & \textbf{Awareness Content} & \textbf{Utility Category} \\
\midrule
ROI calculation & Financial consequence awareness & $A^{self}_{INU,F}$ \\
Ethical reasoning & Moral consequence awareness & $A^{comm}_{IDN}$ \\
Strategic thinking & Social consequence awareness & $A^{self}_{KNU}$ \\
Health evaluation & Long-term physical awareness & $A^{self}_{INU,P,t=3}$ \\
Career planning & Future identity awareness & $A^{self}_{IDN,t>0}$ \\
\bottomrule
\end{tabular}
\end{table}

%------------------------------------------------------------------------------
\subsubsection{System Level Interactions}
%------------------------------------------------------------------------------

\paragraph{The Hyper-Dissonance Rule.}
A critical constraint: Awareness cannot be simultaneously maximal in SYS 0 
and SYS 2:

\begin{equation}
\boxed{A_{SYS0}(t^*) = 1 \quad \Rightarrow \quad A_{SYS2}(t^*) < 1}
\label{eq:hyper-dissonance}
\end{equation}

When visceral reactions dominate, reflection is \emph{structurally suppressed}.

\paragraph{Implication.}
This explains why:
\begin{itemize}
\item Cravings override health knowledge
\item Fear prevents rational evaluation
\item Anger bypasses deliberation
\item Sexual arousal impairs judgment
\end{itemize}

\paragraph{Transition Dynamics.}
Awareness can shift across system levels:

\begin{equation}
SYS0 \rightarrow SYS1 \rightarrow SYS2 \quad \text{(Ascending: reflex → reflection)}
\end{equation}

\begin{equation}
SYS2 \rightarrow SYS1 \rightarrow SYS0 \quad \text{(Descending: deliberation → automaticity)}
\end{equation}

\paragraph{Example: Anger Management.}
\begin{enumerate}
\item SYS 0: Visceral anger response (automatic)
\item SYS 1: Emotional awareness of anger (``I'm furious'')
\item SYS 2: Reflective evaluation (``Should I act on this?'')
\end{enumerate}

Successful anger management requires allowing time for transition from 
SYS 0 to SYS 2 (``count to ten'').

%------------------------------------------------------------------------------
\subsubsection{System Levels and Utility Categories}
%------------------------------------------------------------------------------

Different utility categories are typically processed at different system levels:

\begin{table}[htbp]
\centering
\caption{Typical System Level by Utility Category}
\label{tab:system-utility-mapping}
\renewcommand{\arraystretch}{1.3}
\begin{tabular}{@{}llll@{}}
\toprule
\textbf{Category} & \textbf{Instant (SYS 0)} & \textbf{Affective (SYS 1)} & \textbf{Reflective (SYS 2)} \\
\midrule
$A^{self}_{INU}$ & Pain, hunger, craving & Emotional wellbeing & Financial planning \\
$A^{self}_{KNU}$ & -- & Group emotions & Strategic thinking \\
$A^{self}_{IDN}$ & Public humiliation & Pride, shame & Career identity \\
$A^{comm}_{INU}$ & -- & Empathy, compassion & Ethical reasoning \\
$A^{comm}_{KNU}$ & -- & -- & Abstract collective \\
$A^{comm}_{IDN}$ & -- & Empathic identification & Deep perspective \\
\bottomrule
\end{tabular}
\end{table}

\paragraph{Key Observation.}
$A^{comm}_{KNU}$ (community collective utility) operates \emph{only} at SYS 2.
This is why collective action problems are so difficult---the relevant 
awareness requires effortful, abstract reasoning.

%------------------------------------------------------------------------------

%------------------------------------------------------------------------------
\subsubsection{Cognitive Structure: AWX-A6a}
%------------------------------------------------------------------------------

Reflexive awareness (SYS 2) is not uniformly distributed across individuals. 
AWX-A6a formalizes the cognitive prerequisites:

\begin{axiom}[AWX-A6a: Kognitive Struktur]
Reflexive Awareness depends on literacy, media competence, and logical 
modeling capability.

\textbf{Formal Specification:}
\begin{equation}
A_{refl} = f(\text{Lit}, \text{MedKom}, \text{LogMod})
\end{equation}
where the functional form may be linear or nonlinear.
\end{axiom}

\paragraph{Components.}
\begin{itemize}
    \item \textbf{Literacy (Lit)}: Basic reading and numeracy---foundation for 
          processing complex information
    \item \textbf{Media Competence (MedKom)}: Ability to critically evaluate 
          sources, recognize manipulation, filter misinformation
    \item \textbf{Logical Modeling (LogMod)}: Capacity for abstract reasoning, 
          counterfactual thinking, causal inference
\end{itemize}

\paragraph{Heterogeneity Implication.}
This axiom explains why identical information produces different awareness 
outcomes across population segments. A financially literate individual 
develops SYS 2 awareness about retirement savings more easily than someone 
lacking these prerequisites. This connects to AWX-A74 (Education Effects) 
in the heterogeneity cluster.

\subsubsection{BEATRIX Module Reference}
%------------------------------------------------------------------------------

\begin{table}[htbp]
\centering
\caption{Section 11.3 BEATRIX References}
\label{tab:11-3-beatrix}
\renewcommand{\arraystretch}{1.3}
\begin{tabular}{@{}lll@{}}
\toprule
\textbf{Element} & \textbf{Module} & \textbf{Reference} \\
\midrule
System Levels & BCM2\_05\_AWX\_7240\_01 & systemic\_levels \\
SYS 0 (Instant) & MA\_0.00o & Visceral factors \\
SYS 1 (Affective) & MA\_0.00n & Process over levels \\
SYS 2 (Reflective) & MA\_0.00r & Compliance, ethics \\
Transition Dynamics & MA\_0.00p & Meta-feedback \\
\bottomrule
\end{tabular}
\end{table}

%%%%%%%%%%%%%%%%%%%%%%%%%%%%%%%%%%%%%%%%%%%%%%%%%%%%%%%%%%%%%%%%%%%%%%%%%%%%%%%
% END OF SECTION 11.3
%%%%%%%%%%%%%%%%%%%%%%%%%%%%%%%%%%%%%%%%%%%%%%%%%%%%%%%%%%%%%%%%%%%%%%%%%%%%%%%
